\documentclass[12pt, a4paper]{article}

% --- Packages ---
\usepackage[margin=2.5cm]{geometry}
\usepackage{amsmath, amssymb, amsthm}
\usepackage{mathtools}
\usepackage{bm}
\usepackage{hyperref}
\usepackage{booktabs}
\usepackage{array}
\usepackage{xcolor}
\usepackage{listings}
\usepackage{titlesec}
\usepackage{parskip}
\usepackage{enumitem}
\usepackage{fancyhdr}
\usepackage{caption}
\usepackage{float}
\usepackage{setspace}

% --- Colors ---
\definecolor{codeblue}{RGB}{0, 80, 160}
\definecolor{codegray}{RGB}{245, 245, 245}
\definecolor{commentgreen}{RGB}{0, 128, 0}
\definecolor{keywordred}{RGB}{160, 0, 0}

% --- Code style ---
\lstdefinestyle{pythonstyle}{
    backgroundcolor=\color{codegray},
    commentstyle=\color{commentgreen}\itshape,
    keywordstyle=\color{keywordred}\bfseries,
    numberstyle=\tiny\color{gray},
    stringstyle=\color{codeblue},
    basicstyle=\ttfamily\small,
    breaklines=true,
    captionpos=b,
    numbers=left,
    numbersep=5pt,
    showspaces=false,
    showstringspaces=false,
    language=Python,
    frame=single,
    rulecolor=\color{gray!40},
    xleftmargin=1em,
    framexleftmargin=1em,
}
\lstset{style=pythonstyle}

% --- Header / Footer ---
\pagestyle{fancy}
\fancyhf{}
\setlength{\headheight}{15pt}
\rhead{Black--Scholes PDE Solver}
\lhead{Mathematical Reference}
\cfoot{\thepage}

% --- Theorem environments ---
\newtheorem{definition}{Definition}[section]
\newtheorem{theorem}{Theorem}[section]
\newtheorem{remark}{Remark}[section]
\newtheorem{proposition}{Proposition}[section]

% --- Title ---
\title{
    \vspace{-1cm}
    \textbf{\Large Black--Scholes PDE Solver} \\[0.4em]
    \large Mathematical Documentation \\[0.2em]
    \normalsize A complete exposition of all numerical methods, discretisations,
    stability analyses, and convergence results implemented in the codebase
}
\author{RobinBasso}
\date{\today}

% ============================================================
\begin{document}
% ============================================================

\maketitle
\thispagestyle{fancy}

\begin{abstract}
This document provides a rigorous mathematical account of every computation
performed in the \texttt{black-scholes-pde} project. The codebase implements
three finite difference schemes for the Black--Scholes partial differential
equation (PDE): an explicit (forward Euler in time) scheme, a
Crank--Nicolson scheme, and Crank--Nicolson with Rannacher smoothing. All
methods are validated against the closed-form Black--Scholes formula and
subjected to systematic spatial and temporal convergence studies. Every
formula, coefficient, matrix, stability condition, and convergence result
stated in this document corresponds directly to the source code in
\texttt{src/solver.py}, \texttt{src/utils.py}, and \texttt{tests/convergence.py}.
\end{abstract}

\tableofcontents
\newpage

% ============================================================
\section{Notation and Conventions}
% ============================================================

\begin{center}
\begin{tabular}{cl}
\toprule
\textbf{Symbol} & \textbf{Meaning} \\
\midrule
$S$ & Underlying asset price \\
$t$ & Time (calendar time, $t \in [0, T]$) \\
$\tau = T - t$ & Time to maturity (backward time) \\
$V(S, t)$ & Option price as a function of $S$ and $t$ \\
$K$ & Strike price (default: 100) \\
$T$ & Maturity (default: 1.0 year) \\
$r$ & Risk-free interest rate (default: 0.05) \\
$\sigma$ & Volatility (default: 0.20) \\
$S_{\max}$ & Truncation of the spatial domain (default: 200) \\
$M$ & Number of spatial intervals ($M+1$ grid points) \\
$N$ & Number of time steps \\
$\Delta S = S_{\max}/M$ & Spatial step size \\
$\Delta t = T/N$ & Time step size \\
$S_i = i\,\Delta S$ & $i$-th spatial grid point, $i = 0, 1, \ldots, M$ \\
$t^n = n\,\Delta t$ & $n$-th time level (forward), $n = 0, 1, \ldots, N$ \\
$V_i^n$ & Numerical approximation to $V(S_i, t^n)$ \\
\bottomrule
\end{tabular}
\end{center}

\medskip
\noindent
The code solves the PDE \textbf{backwards in time} from $t = T$ (terminal
condition) to $t = 0$ (option price today), storing the solution on a grid
\texttt{grid[i, j]} where the column index $j$ runs from $0$ (today) to
$N$ (maturity). The analytical solution is stored at column $j = 0$ after
backward stepping.

% ============================================================
\section{The Black--Scholes PDE}
% ============================================================

\subsection{Model and Equation}

Under the Black--Scholes model, the no-arbitrage price $V(S, t)$ of a
European option satisfies the PDE:
\begin{equation}
    \boxed{
    \frac{\partial V}{\partial t}
    + \frac{1}{2}\sigma^2 S^2 \frac{\partial^2 V}{\partial S^2}
    + r S \frac{\partial V}{\partial S}
    - r V = 0,
    \qquad (S, t) \in (0, S_{\max}) \times [0, T).
    }
    \label{eq:bspde}
\end{equation}
The three terms beyond $\partial V/\partial t$ have clear financial
interpretations:
\begin{itemize}[noitemsep]
    \item $\frac{1}{2}\sigma^2 S^2 \partial^2 V/\partial S^2$: convexity / gamma term (diffusion).
    \item $r S\, \partial V/\partial S$: delta term (advection).
    \item $-rV$: discounting (reaction / zeroth-order term).
\end{itemize}

\subsection{Terminal Condition}

For a European call option with strike $K$:
\begin{equation}
    V(S, T) = \max(S - K,\; 0).
    \label{eq:terminal}
\end{equation}
This is the \textbf{payoff function}. In the code it is initialised as:
\begin{lstlisting}
grid[:, -1] = np.maximum(S - K, 0)
\end{lstlisting}
where the last column (\texttt{j = N}) corresponds to $t = T$.

\subsection{Boundary Conditions}
\label{sec:bc}

\paragraph{Lower boundary ($S = 0$).}
When the asset price is zero, it stays zero under the geometric Brownian
motion model, so the call option is worthless:
\begin{equation}
    V(0, t) = 0 \qquad \forall\, t \in [0, T].
    \label{eq:bc_lower}
\end{equation}

\paragraph{Upper boundary ($S = S_{\max}$).}
For large $S$, the probability of the option expiring out-of-the-money
approaches zero and the option behaves like a forward contract. The standard
approximation is:
\begin{equation}
    V(S_{\max}, t) \approx S_{\max} - K e^{-r(T - t)}.
    \label{eq:bc_upper}
\end{equation}
In the code, this is precomputed over all time levels:
\begin{lstlisting}
time = np.linspace(0, T, N + 1)
grid[-1, :] = S_max - K * np.exp(-r * (T - time))
\end{lstlisting}

% ============================================================
\section{Analytical Solution (\texttt{utils.py})}
% ============================================================

\subsection{Black--Scholes Closed-Form Formula}

The exact price of a European call under the Black--Scholes model is:
\begin{equation}
    \boxed{V(S, t) = S\,\Phi(d_1) - K e^{-r(T-t)}\,\Phi(d_2)}
    \label{eq:bs_formula}
\end{equation}
where $\Phi$ is the standard normal CDF and:
\begin{align}
    d_1 &= \frac{\ln(S/K) + (r + \tfrac{1}{2}\sigma^2)(T-t)}{\sigma\sqrt{T-t}},
    \label{eq:d1} \\
    d_2 &= d_1 - \sigma\sqrt{T-t}.
    \label{eq:d2}
\end{align}
The implementation in \texttt{black\_scholes\_call} evaluates this at $t = 0$
(i.e., $T - t = T$), masking out $S = 0$ to avoid $\ln(0)$:
\begin{lstlisting}
mask   = S > 0
d1     = (np.log(S_pos / K) + (r + 0.5*sigma**2)*T) / (sigma*np.sqrt(T))
d2     = d1 - sigma * np.sqrt(T)
price[mask] = S_pos * norm.cdf(d1) - K * np.exp(-r*T) * norm.cdf(d2)
\end{lstlisting}
This serves as the \textbf{ground truth} for all error computations in the
project.

% ============================================================
\section{Spatial Discretisation}
% ============================================================

\subsection{Grid Construction}

The spatial domain $[0, S_{\max}]$ is divided into $M$ equal intervals:
\[
    S_i = i\,\Delta S, \quad \Delta S = \frac{S_{\max}}{M}, \quad i = 0, 1, \ldots, M.
\]
The grid has $M + 1$ points total. Interior points are $i = 1, \ldots, M-1$;
boundary points are $i = 0$ and $i = M$.

\subsection{Finite Difference Operators}

The spatial derivatives in equation~\eqref{eq:bspde} are approximated by
\textbf{second-order central differences}:
\begin{align}
    \frac{\partial V}{\partial S}\bigg|_{S=S_i}
    &\approx \frac{V_{i+1} - V_{i-1}}{2\,\Delta S},
    \label{eq:fd_first} \\
    \frac{\partial^2 V}{\partial S^2}\bigg|_{S=S_i}
    &\approx \frac{V_{i+1} - 2V_i + V_{i-1}}{\Delta S^2}.
    \label{eq:fd_second}
\end{align}
Both are $O(\Delta S^2)$ accurate, which underpins the second-order spatial
convergence observed in all three schemes.

\subsection{Substitution into the PDE}

Substituting~\eqref{eq:fd_first}--\eqref{eq:fd_second} into~\eqref{eq:bspde},
noting that $S_i = i\,\Delta S$ so $S_i/\Delta S = i$, and using
$S_i^2/\Delta S^2 = i^2$, yields:
\begin{equation}
    \frac{\partial V_i}{\partial t}
    + \frac{\sigma^2 i^2}{2}\bigl(V_{i+1} - 2V_i + V_{i-1}\bigr)
    + \frac{r i}{2}\bigl(V_{i+1} - V_{i-1}\bigr)
    - r V_i = 0.
    \label{eq:semidiscrete}
\end{equation}
This is a \emph{semi-discrete} system (continuous in time, discrete in space),
from which the fully discrete schemes are derived by approximating
$\partial V_i/\partial t$.

% ============================================================
\section{Explicit Finite Difference Scheme (\texttt{solve\_explicit})}
% ============================================================

\subsection{Time Discretisation}

The time derivative is approximated by a \textbf{forward (explicit) Euler}
step — but since we march \emph{backwards} from $t = T$ to $t = 0$, the
update reads:
\[
    \frac{\partial V_i}{\partial t} \approx \frac{V_i^{n+1} - V_i^n}{\Delta t}
    \quad \Longrightarrow \quad
    V_i^n = V_i^{n+1} + \Delta t \left[\text{spatial terms at } n+1\right].
\]
Substituting into~\eqref{eq:semidiscrete} and rearranging:

\begin{equation}
    \boxed{V_i^n = a_i\,V_{i-1}^{n+1} + b_i\,V_i^{n+1} + c_i\,V_{i+1}^{n+1}}
    \label{eq:explicit_update}
\end{equation}

\subsection{Coefficients}

\begin{align}
    a_i &= \frac{1}{2}\,\Delta t\!\left(\sigma^2 i^2 - r i\right),
    \label{eq:ai} \\
    b_i &= 1 - \Delta t\!\left(\sigma^2 i^2 + r\right),
    \label{eq:bi} \\
    c_i &= \frac{1}{2}\,\Delta t\!\left(\sigma^2 i^2 + r i\right).
    \label{eq:ci}
\end{align}

These match exactly the code:
\begin{lstlisting}
a = 0.5 * dt * (sigma**2 * i**2 - r * i)
b = 1 - dt * (sigma**2 * i**2 + r)
c = 0.5 * dt * (sigma**2 * i**2 + r * i)
grid[i, j] = a * grid[i-1, j+1] + b * grid[i, j+1] + c * grid[i+1, j+1]
\end{lstlisting}

Note that $a_i + b_i + c_i = 1 - \Delta t\cdot r$, which is less than 1 for
$r > 0$, consistent with discounting.

\subsection{Stability Analysis}
\label{sec:explicit_stability}

The explicit scheme is \textbf{conditionally stable}. The stability condition
requires that all coefficients $a_i$, $b_i$, $c_i$ be non-negative (a
sufficient condition for monotonicity / maximum principle):
\[
    b_i \geq 0 \quad \Longleftrightarrow \quad
    \Delta t \leq \frac{1}{\sigma^2 i^2 + r}.
\]
The most restrictive constraint is at $i = M$ (largest $S$):
\begin{equation}
    \boxed{\Delta t \lesssim \frac{1}{\sigma^2 M^2}}
    \label{eq:cfl}
\end{equation}
This is a \textbf{CFL-type condition}. In the convergence study, the explicit
scheme is given $N = 2M^2$ time steps to ensure stability:
\begin{lstlisting}
N_explicit = M * M * 2   # ensures dt << 1 / (sigma^2 * M^2)
\end{lstlisting}
Violation of~\eqref{eq:cfl} causes the scheme to diverge, as demonstrated in
the time convergence tests.

\subsection{Truncation Error}

The local truncation error (LTE) of the explicit scheme is:
\[
    \tau = O(\Delta t) + O(\Delta S^2).
\]
Hence the scheme is \textbf{first-order in time} and \textbf{second-order in
space}. When the CFL condition is tight ($\Delta t \sim 1/(\sigma^2 M^2) \sim
\Delta S^2$), both errors are of comparable magnitude.

% ============================================================
\section{Crank--Nicolson Scheme (\texttt{solve\_crank\_nicolson})}
% ============================================================

\subsection{Time Discretisation}

The Crank--Nicolson method averages the spatial operator at time levels
$n$ and $n+1$:
\begin{equation}
    \frac{V_i^n - V_i^{n+1}}{\Delta t}
    = \frac{1}{2}\mathcal{L}_i[V^n]
    + \frac{1}{2}\mathcal{L}_i[V^{n+1}],
    \label{eq:cn_abstract}
\end{equation}
where the spatial operator (from equation~\eqref{eq:semidiscrete}) is:
\[
    \mathcal{L}_i[V]
    = \frac{\sigma^2 i^2}{2}(V_{i+1} - 2V_i + V_{i-1})
    + \frac{ri}{2}(V_{i+1} - V_{i-1}) - rV_i.
\]

\subsection{Coefficient Derivation}

Define half-step coefficients:
\begin{align}
    \alpha_i &= \frac{1}{4}\,\Delta t\!\left(\sigma^2 i^2 - ri\right),
    \label{eq:alpha} \\
    \beta_i  &= -\frac{1}{2}\,\Delta t\!\left(\sigma^2 i^2 + r\right),
    \label{eq:beta} \\
    \gamma_i &= \frac{1}{4}\,\Delta t\!\left(\sigma^2 i^2 + ri\right).
    \label{eq:gamma}
\end{align}
Note that $\alpha_i = a_i/2$ and $\gamma_i = c_i/2$ compared to the explicit
coefficients, and $\beta_i = (b_i - 1)/2$. These match the code:
\begin{lstlisting}
alpha = 0.25 * dt * (sigma**2 * i**2 - r * i)
beta  = -0.5 * dt * (sigma**2 * i**2 + r)
gamma = 0.25 * dt * (sigma**2 * i**2 + r * i)
\end{lstlisting}

\subsection{Linear System}

Rearranging~\eqref{eq:cn_abstract} moves the unknowns ($V^n$) to the left and
the knowns ($V^{n+1}$) to the right, giving a tridiagonal linear system for
the interior points $i = 1, \ldots, M-1$:
\begin{equation}
    \mathbf{A}\,\bm{V}^n = \mathbf{B}\,\bm{V}^{n+1} + \bm{b}_{\text{bc}},
    \label{eq:linsys}
\end{equation}
where $\bm{V}^n = (V_1^n, V_2^n, \ldots, V_{M-1}^n)^\top$ and the
$(M-1)\times(M-1)$ tridiagonal matrices are:

\paragraph{Matrix $\mathbf{A}$ (implicit side):}
\begin{equation}
    \mathbf{A} =
    \begin{pmatrix}
    1-\beta_1   & -\gamma_1  & 0          & \cdots & 0 \\
    -\alpha_2   & 1-\beta_2  & -\gamma_2  & \cdots & 0 \\
    0           & -\alpha_3  & 1-\beta_3  & \cdots & 0 \\
    \vdots      & \ddots     & \ddots     & \ddots & \vdots \\
    0           & \cdots     & 0          & -\alpha_{M-1} & 1-\beta_{M-1}
    \end{pmatrix}
    \label{eq:A_matrix}
\end{equation}

\paragraph{Matrix $\mathbf{B}$ (explicit side):}
\begin{equation}
    \mathbf{B} =
    \begin{pmatrix}
    1+\beta_1   & \gamma_1   & 0          & \cdots & 0 \\
    \alpha_2    & 1+\beta_2  & \gamma_2   & \cdots & 0 \\
    0           & \alpha_3   & 1+\beta_3  & \cdots & 0 \\
    \vdots      & \ddots     & \ddots     & \ddots & \vdots \\
    0           & \cdots     & 0          & \alpha_{M-1} & 1+\beta_{M-1}
    \end{pmatrix}
    \label{eq:B_matrix}
\end{equation}

Note that $\mathbf{A}$ and $\mathbf{B}$ satisfy $\mathbf{A} + \mathbf{B} =
2\mathbf{I}$, up to boundary corrections. In the code:
\begin{lstlisting}
A[row, row]   = 1 - beta;  B[row, row]   = 1 + beta
A[row, row-1] = -alpha;    B[row, row-1] = alpha
A[row, row+1] = -gamma;    B[row, row+1] = gamma
\end{lstlisting}

\subsection{Boundary Corrections}

The right-hand side vector $\bm{b}_{\text{bc}}$ accounts for the
boundary values at $i = 0$ and $i = M$ which appear in the stencil for
the first and last interior rows. For row $i = 1$ (lower boundary):
\[
    [\bm{b}_{\text{bc}}]_1
    = \alpha_1 \bigl(V_0^n + V_0^{n+1}\bigr)
    = \alpha_1 \cdot 0 = 0,
\]
but since $V_0^n = 0$ always, the correction is zero for the lower boundary.
For row $i = M-1$ (upper boundary):
\[
    [\bm{b}_{\text{bc}}]_{M-1}
    = \gamma_{M-1}\bigl(V_M^n + V_M^{n+1}\bigr).
\]
This is added in the time loop:
\begin{lstlisting}
# Lower boundary (i=1): V_0 = 0, so correction is zero
alpha_1 = 0.25 * dt * (sigma**2 * 1**2 - r * 1)
rhs[0] += alpha_1 * (grid[0, j] + grid[0, j+1])

# Upper boundary (i=M-1)
gamma_Mm1 = 0.25 * dt * (sigma**2 * (M-1)**2 + r * (M-1))
rhs[-1]  += gamma_Mm1 * (grid[M, j] + grid[M, j+1])
\end{lstlisting}

\subsection{Linear Solve}

At each time step, the system $\mathbf{A}\bm{x} = \text{rhs}$ is solved using
\texttt{numpy.linalg.solve}, which applies LU decomposition. Since $\mathbf{A}$
is a fixed tridiagonal matrix (time-independent coefficients), it could be
pre-factorised for efficiency; the current code recomputes the solve at each
step.

\subsection{Truncation Error and Stability}

The Crank--Nicolson scheme has:
\begin{itemize}[noitemsep]
    \item \textbf{Truncation error}: $O(\Delta t^2) + O(\Delta S^2)$ — second-order in both.
    \item \textbf{Stability}: \textbf{unconditionally stable} (for the diffusion-dominated case), because $\mathbf{A}$ is strictly diagonally dominant for $\beta_i < 0$, which holds whenever $r > 0$.
\end{itemize}

The unconditional stability means $N$ can be set equal to $M$ (rather than
$M^2$) in the convergence study:
\begin{lstlisting}
N_cn = M   # No CFL constraint
\end{lstlisting}

% ============================================================
\section{Rannacher Smoothing (\texttt{solve\_rannacher})}
% ============================================================

\subsection{Motivation: Payoff Regularity}

The call payoff $\max(S - K, 0)$ is \textbf{continuous but not differentiable}
at $S = K$:
\[
    \frac{\partial V}{\partial S}\bigg|_{S=K^+} = 1,
    \qquad
    \frac{\partial V}{\partial S}\bigg|_{S=K^-} = 0.
\]
This kink in the initial condition degrades the temporal convergence of
Crank--Nicolson from theoretical $O(\Delta t^2)$ to approximately $O(\Delta t)$
in practice, because the averaging of the operator at two time levels amplifies
the discontinuity in the derivative.

\subsection{The Rannacher Procedure}

The Rannacher method (1984) applies \textbf{fully implicit (backward) Euler}
for the first $n_{\text{smooth}}$ time steps (default: 4) immediately after
the terminal condition, before switching to Crank--Nicolson:
\begin{equation}
    \text{Method at step } j =
    \begin{cases}
        \text{Implicit Euler} & \text{if } j \geq N - n_{\text{smooth}} \\
        \text{Crank--Nicolson} & \text{otherwise}
    \end{cases}
\end{equation}
The implicit Euler steps damp the oscillations caused by the non-smooth
payoff; subsequent Crank--Nicolson steps then inherit a smooth solution.

\subsection{Implicit Euler Coefficients}

The implicit Euler discretisation is:
\[
    \frac{V_i^n - V_i^{n+1}}{\Delta t} = \mathcal{L}_i[V^n],
\]
which gives the linear system $\mathbf{A}_{\text{imp}}\,\bm{V}^n =
\bm{V}^{n+1}$ with tridiagonal matrix coefficients:
\begin{align}
    \alpha_i^{\text{(IE)}} &= \frac{1}{2}\,\Delta t(\sigma^2 i^2 - ri),
    \label{eq:alpha_imp} \\
    \beta_i^{\text{(IE)}}  &= -\Delta t(\sigma^2 i^2 + r),
    \label{eq:beta_imp} \\
    \gamma_i^{\text{(IE)}} &= \frac{1}{2}\,\Delta t(\sigma^2 i^2 + ri).
    \label{eq:gamma_imp}
\end{align}
Note that $\alpha_i^{\text{(IE)}} = 2\alpha_i$ and $\gamma_i^{\text{(IE)}} =
2\gamma_i$ — exactly twice the Crank--Nicolson half-coefficients. This matches
the code:
\begin{lstlisting}
alpha_e = 0.5 * dt * (sigma**2 * i**2 - r * i)   # = 2 * alpha_CN
beta_e  = -dt  * (sigma**2 * i**2 + r)             # = 2 * beta_CN
gamma_e = 0.5 * dt * (sigma**2 * i**2 + r * i)    # = 2 * gamma_CN

A_imp[row, row-1] = -alpha_e
A_imp[row, row]   = 1 - beta_e
A_imp[row, row+1] = -gamma_e
\end{lstlisting}

\subsection{Right-Hand Side During Implicit Euler Phase}

During the Rannacher smoothing phase ($j \geq N - n_{\text{smooth}}$), the
RHS for the implicit Euler step is simply $\bm{V}^{n+1}$ with boundary
corrections from both boundaries:
\begin{lstlisting}
rhs = grid[1:M, j+1].copy()
rhs[0]  += 0.5*dt*(sigma**2*1**2 - r*1)    * (grid[0, j] + grid[0, j+1])
rhs[-1] += 0.5*dt*(sigma**2*(M-1)**2 + r*(M-1)) * (grid[M, j] + grid[M, j+1])
\end{lstlisting}

\subsection{Truncation Error}

\begin{itemize}[noitemsep]
    \item Implicit Euler: $O(\Delta t) + O(\Delta S^2)$ — first-order in time.
    \item Crank--Nicolson: $O(\Delta t^2) + O(\Delta S^2)$ — second-order in time (theoretical).
\end{itemize}
With the Rannacher procedure, the observed convergence is
\textbf{approximately first-order in time} rather than second-order, because
the $O(\Delta t)$ error from the initial smoothing steps dominates. However,
the smoothing eliminates the spurious oscillations that would otherwise
completely destroy the convergence rate of plain Crank--Nicolson on
non-smooth payoffs.

% ============================================================
\section{Convergence Analysis (\texttt{tests/convergence.py})}
% ============================================================

\subsection{Error Metric}

All convergence studies use the \textbf{maximum absolute error} (sup-norm)
over the spatial grid at $t = 0$:
\begin{equation}
    \mathcal{E} = \max_{i=0,\ldots,M} \left|V_i^{\text{num}} - V^{\text{exact}}(S_i)\right|,
    \label{eq:error}
\end{equation}
implemented as:
\begin{lstlisting}
error = np.max(np.abs(price_numerical - price_exact))
\end{lstlisting}

\subsection{Spatial Convergence Study}

\paragraph{Setup.} Spatial resolution is varied: $M \in \{25, 50, 100, 200\}$.
For each $M$:
\begin{itemize}[noitemsep]
    \item Crank--Nicolson: $N = M$ (time error negligible, dominated by space).
    \item Explicit: $N = 2M^2$ (ensures CFL stability without time error dominating).
\end{itemize}

\paragraph{Expected result.} Both schemes use second-order central differences
in space, giving:
\[
    \mathcal{E} = O(\Delta S^2) = O(M^{-2}).
\]

\paragraph{Observed result.} Both schemes track the $O(M^{-2})$ reference
slope on a log-log plot. However, neither exactly achieves the theoretical
rate: the kink in the payoff at $S = K$ locally reduces the spatial regularity
and slightly degrades the observed order.

\subsection{Time Convergence Study}

\paragraph{Setup.} Fixed $M = 50$ (spatial error fixed). The stability limit
is:
\[
    \Delta t_{\text{stab}} = \frac{1}{\sigma^2 M^2} = \frac{1}{0.04 \times 2500}
    = \frac{1}{100} = 0.01,
    \qquad
    N_{\min} = \left\lfloor\frac{T}{\Delta t_{\text{stab}}}\right\rfloor = 100.
\]
Both schemes are run at $N \in \{N_{\min}, 2N_{\min}, 4N_{\min}, 8N_{\min},
16N_{\min}\}$.

\paragraph{Explicit scheme instability.} When $N < N_{\min}$, the explicit
scheme violates the CFL condition~\eqref{eq:cfl} and diverges. The convergence
study shows this by using $N$ values just above the stability boundary.

\paragraph{Saturation phenomenon.} Once the explicit scheme is stabilised,
increasing $N$ barely reduces the error, because the dominant error is now
\emph{spatial} (fixed $M$). The time error $O(\Delta t) = O(1/N)$ quickly
becomes negligible compared to $O(\Delta S^2) = O(M^{-2})$.

\paragraph{Crank--Nicolson.} Similar saturation occurs: the scheme is already
accurate enough in time for reasonable $N$, and the spatial error floor is
reached quickly.

\subsection{Time Convergence with Rannacher Smoothing}

\paragraph{Setup.} Fixed $M = 200$ (fine spatial grid to suppress spatial
error). Time steps $N \in \{100, 200, 400, 800, 1600\}$.

\paragraph{Observed convergence.} The Rannacher scheme achieves approximately
$O(\Delta t) = O(N^{-1})$ in time, matching the $O(dt)$ reference slope on the
log-log plot. This is \textbf{better} than plain Crank--Nicolson on the same
non-smooth payoff (which stagnates or oscillates), but \textbf{worse} than
theoretical second-order because the implicit Euler smoothing steps introduce
$O(\Delta t)$ error.

\begin{remark}
    The $O(\Delta t)$ convergence of Rannacher is not a failure of the method.
    It is the \emph{correct} asymptotic rate given the payoff regularity: the
    gain over plain Crank--Nicolson lies in the \emph{constant} being much
    smaller and in the elimination of oscillatory artefacts.
\end{remark}

% ============================================================
\section{End-to-End Computational Summary}
% ============================================================

The full computation for a single numerical solution can be written as:

\bigskip
\noindent\textbf{Step 1 — Grid initialisation.}
\[
    S_i = i\,\frac{S_{\max}}{M}, \quad i = 0,\ldots,M; \qquad
    t^n = n\,\frac{T}{N}, \quad n = 0,\ldots,N.
\]

\noindent\textbf{Step 2 — Terminal condition.}
\[
    V_i^N = \max(S_i - K,\; 0), \quad i = 0,\ldots,M.
\]

\noindent\textbf{Step 3 — Boundary conditions (all $n$).}
\[
    V_0^n = 0, \qquad V_M^n = S_{\max} - K e^{-r(T - t^n)}.
\]

\noindent\textbf{Step 4 — Backward time stepping} (for $j = N-1, N-2, \ldots, 0$):
\begin{itemize}[noitemsep]
    \item \textbf{Explicit}: $V_i^j = a_i V_{i-1}^{j+1} + b_i V_i^{j+1} + c_i V_{i+1}^{j+1}$, $i = 1,\ldots,M-1$.
    \item \textbf{Crank--Nicolson}: solve $\mathbf{A}\bm{V}^j = \mathbf{B}\bm{V}^{j+1} + \bm{b}_{\text{bc}}$.
    \item \textbf{Rannacher}: use implicit Euler for $j \geq N - n_{\text{smooth}}$, then Crank--Nicolson.
\end{itemize}

\noindent\textbf{Step 5 — Output.}
\[
    \text{Return } (S_0, S_1, \ldots, S_M) \text{ and } (V_0^0, V_1^0, \ldots, V_M^0).
\]

% ============================================================
\section{Scheme Comparison}
% ============================================================

\begin{table}[H]
\centering
\caption{Summary comparison of the three implemented schemes}
\begin{tabular}{llllll}
\toprule
\textbf{Scheme} & \textbf{Time order} & \textbf{Space order} & \textbf{Stability} & \textbf{System} & \textbf{$N$ needed} \\
\midrule
Explicit        & $O(\Delta t)$       & $O(\Delta S^2)$      & Conditional (CFL)  & Explicit        & $O(M^2)$ \\
Crank--Nicolson & $O(\Delta t^2)$     & $O(\Delta S^2)$      & Unconditional      & Tridiagonal     & $O(M)$ \\
Rannacher       & $O(\Delta t)$*      & $O(\Delta S^2)$      & Unconditional      & Tridiagonal     & $O(M)$ \\
\bottomrule
\end{tabular}
\caption*{*Practical rate on non-smooth payoff; theoretical CN rate is $O(\Delta t^2)$.}
\end{table}

% ============================================================
\section{Module Architecture Summary}
% ============================================================

\begin{table}[H]
\centering
\caption{Modules and their mathematical roles}
\begin{tabular}{llp{8.5cm}}
\toprule
\textbf{File} & \textbf{Function} & \textbf{Mathematical Content} \\
\midrule
\texttt{utils.py} & \texttt{black\_scholes\_call} & Closed-form formula (equations~\ref{eq:bs_formula}--\ref{eq:d2}) \\
\texttt{solver.py} & \texttt{solve\_explicit} & Explicit update (equation~\ref{eq:explicit_update}), coefficients~\eqref{eq:ai}--\eqref{eq:ci} \\
& \texttt{solve\_crank\_nicolson} & Tridiagonal system~\eqref{eq:linsys}, matrices~\eqref{eq:A_matrix}--\eqref{eq:B_matrix} \\
& \texttt{solve\_rannacher} & Hybrid IE + CN, coefficients~\eqref{eq:alpha_imp}--\eqref{eq:gamma_imp} \\
\texttt{convergence.py} & \texttt{convergence\_study} & Spatial $\mathcal{E}$ vs.\ $M$ (equation~\ref{eq:error}) \\
& \texttt{time\_convergence} & Temporal $\mathcal{E}$ vs.\ $N$, CFL illustration \\
& \texttt{time\_convergence\_rannacher} & Rannacher $\mathcal{E}$ vs.\ $N$, $O(\Delta t)$ rate \\
\texttt{test\_solver.py} & \texttt{run\_test} & Point-wise comparison of all three methods against analytical \\
\bottomrule
\end{tabular}
\end{table}

% ============================================================
\section{Limitations and Potential Extensions}
% ============================================================

The current implementation is a clean, instructive codebase with a few
practical limitations worth noting for production use:

\begin{enumerate}
    \item \textbf{European call only.} The terminal condition and upper
          boundary are hardcoded for a vanilla call. American options require
          a free-boundary (complementarity) formulation; exotic payoffs require
          modified terminal conditions.

    \item \textbf{Non-uniform grids.} The uniform $\Delta S$ grid concentrates
          grid points uniformly across $[0, S_{\max}]$. A sinh-stretching or
          log-transform grid would place more points near $S = K$, improving
          accuracy at the kink and potentially restoring second-order spatial
          convergence in the convergence constant.

    \item \textbf{Direct linear solver.} \texttt{np.linalg.solve} uses
          full LU decomposition, $O(M^3)$ per time step. For tridiagonal
          systems, the Thomas algorithm is $O(M)$; using
          \texttt{scipy.linalg.solve\_banded} or \texttt{scipy.sparse} would
          give a significant speedup for large $M$.

    \item \textbf{Constant $\mathbf{A}$.} The matrix $\mathbf{A}$ is
          time-independent (coefficients do not depend on $t^n$). It is
          rebuilt and refactorised at each time step in the current code; a
          single pre-factorisation would reduce computational cost.

    \item \textbf{Single asset, single factor.} The code solves a
          one-dimensional PDE. Multi-asset options (basket, spread) require
          higher-dimensional PDE solvers (ADI methods, Monte Carlo).

    \item \textbf{Fixed parameters.} $\sigma$ and $r$ are constant. Local
          volatility models (Dupire) or stochastic volatility models (Heston)
          introduce $\sigma = \sigma(S, t)$ or additional state variables.
\end{enumerate}

Natural extensions include: Richardson extrapolation for higher-order time
accuracy without Rannacher, PSOR (projected SOR) for American options, and
a log-space transformation $x = \ln(S/K)$ which converts the PDE to
constant-coefficient form and simplifies grid design.

% ============================================================
\section*{References}
% ============================================================

\begin{itemize}[noitemsep]
    \item Black, F. \& Scholes, M. (1973). The pricing of options and
          corporate liabilities. \textit{Journal of Political Economy}, 81(3),
          637--654.
    \item Rannacher, R. (1984). Finite element solution of diffusion problems
          with irregular data. \textit{Numerische Mathematik}, 43(2), 309--327.
    \item Wilmott, P., Dewynne, J. \& Howison, S. (1993). \textit{Option
          Pricing: Mathematical Models and Computation}. Oxford Financial
          Press.
    \item Strikwerda, J. C. (2004). \textit{Finite Difference Schemes and
          Partial Differential Equations} (2nd ed.). SIAM.
    \item Seydel, R. U. (2009). \textit{Tools for Computational Finance}
          (4th ed.). Springer.
\end{itemize}

\end{document}
